\chapter{Movimento e Direção}

\begin{objetivos}
\begin{itemize}
\item Usar os blocos de \textbf{Movimento} para mexer os atores;
\item Entender o plano do palco com as coordenadas \textbf{x} e \textbf{y};
\item Fazer um ator girar, deslizar e repetir um movimento.
\end{itemize}
\end{objetivos}

\section{Os blocos de movimento}

Os blocos azuis \bloco[scMotion]{Movimento} dizem para onde o ator vai. Os mais usados são:

\begin{center}
\renewcommand{\arraystretch}{1.8}
\begin{tabular}{@{}ll@{}}
\bloco[scMotion]{mova \entrada{10} passos} & anda para a frente \\
\bloco[scMotion]{gire \entrada{15} graus} & vira o ator \\
\bloco[scMotion]{vá para x: \entrada{0} y: \entrada{0}} & vai direto para um ponto \\
\bloco[scMotion]{deslize por \entrada{1} seg até x: \entrada{0} y: \entrada{0}} & vai deslizando até lá \\
\end{tabular}
\end{center}

\begin{dica}
Passos \textbf{negativos} fazem o ator andar para trás! Experimente \bloco[scMotion]{mova \entrada{$-10$} passos}.
\end{dica}

\section{O plano do palco: x e y}

O palco é como um \textbf{papel quadriculado}. Cada lugar tem um endereço com dois números: o \textbf{x} (esquerda--direita) e o \textbf{y} (embaixo--em cima). O centro é o ponto \textbf{(0, 0)}.

\begin{center}
\begin{tikzpicture}[font=\sffamily\small]
  \draw[fill=white,draw=ink,line width=1.3pt,rounded corners=4pt] (-4.8,-3.6) rectangle (4.8,3.6);
  \draw[gridc!70!ink,line width=.4pt] (-4.8,-3.6) grid[step=.6] (4.8,3.6);
  \draw[ink,-{Stealth},line width=1.1pt] (-4.8,0)--(5.1,0) node[right]{\bfseries x};
  \draw[ink,-{Stealth},line width=1.1pt] (0,-3.6)--(0,3.9) node[above]{\bfseries y};
  \node[fill=scMotion,draw=ink,text=white,inner sep=2pt,rounded corners=2pt,font=\sffamily\bfseries\scriptsize] at (-4.4,-.3) {$-240$};
  \node[fill=scMotion,draw=ink,text=white,inner sep=2pt,rounded corners=2pt,font=\sffamily\bfseries\scriptsize] at (4.4,-.3) {$240$};
  \node[fill=scOperators,draw=ink,text=white,inner sep=2pt,rounded corners=2pt,font=\sffamily\bfseries\scriptsize] at (.6,3.3) {$180$};
  \node[fill=scOperators,draw=ink,text=white,inner sep=2pt,rounded corners=2pt,font=\sffamily\bfseries\scriptsize] at (.65,-3.3) {$-180$};
  \fill[scMyBlocks,draw=ink] (2.4,1.8) circle(.13) node[above right,font=\sffamily\bfseries\footnotesize]{(120, 90)};
  \fill[scLooks,draw=ink] (-2.4,-1.2) circle(.13) node[below left,font=\sffamily\bfseries\footnotesize]{($-120$, $-60$)};
  \fill[scEvents,draw=ink] (0,0) circle(.13);
\end{tikzpicture}
\end{center}

\begin{definicao}[Resumo]
\textbf{x} mede \emph{para o lado}: positivo é direita, negativo é esquerda.\\
\textbf{y} mede \emph{para cima}: positivo é para cima, negativo é para baixo.
\end{definicao}

\section{Girando}

O bloco \bloco[scMotion]{gire \entrada{15} graus} vira o ator para a direita. Uma volta completa tem \textbf{360 graus}. Girar 90 graus é um quarto de volta, como a curva de uma esquina.

\begin{missao}[Missão 2: Dança do gatinho]
Faça o gatinho andar, girar, andar de novo e deslizar de volta para o centro do palco.
\end{missao}

\begin{codigo}[Uma solução possível]
\begin{pilha}
\sbh[scEvents]{quando \bandeira\ for clicado}
\bk[scMotion]{mova \entrada{50} passos}
\bk[scMotion]{gire \entrada{90} graus}
\bk[scMotion]{mova \entrada{50} passos}
\bk[scMotion]{deslize por \entrada{1} seg até x: \entrada{0} y: \entrada{0}}
\end{pilha}
\end{codigo}

\section{Repetindo movimentos}

Imagine ter que arrastar o mesmo bloco 4 vezes! Para isso existe o bloco laranja \bloco[scControl]{repita \entrada{4} vezes}. Tudo o que estiver dentro dele é repetido.

\begin{codigo}[Desenhando um quadrado]
\begin{pilha}
\sbh[scEvents]{quando \bandeira\ for clicado}
\sbc[scControl]{repita \entrada{4} vezes}
\bk[scMotion]{mova \entrada{100} passos}
\bk[scMotion]{gire \entrada{90} graus}
\sbcend
\end{pilha}
\end{codigo}

\begin{vocabulario}[Vocabulário: repetição]
Repetir um conjunto de ordens várias vezes se chama \textbf{laço} (em inglês, \textit{loop}). Ele evita que a gente copie blocos sem parar.
\end{vocabulario}

\begin{pensaassim}[Pense assim: escovar os dentes]
Ninguém diz ``mova a escova, mova a escova, mova a escova\ldots'' trinta vezes. Dizemos só: \emph{``escove por 30 movimentos''}. O bloco \bloco[scControl]{repita} faz isso: você escreve o passo \textbf{uma vez} e diz \textbf{quantas vezes} repetir. No quadrado acima, \emph{mova 100 passos e gire 90 graus} é um ``lado'', e o gatinho repete 4 lados.
\end{pensaassim}

\section{Hora de praticar}

\begin{exercicio}[Vá para o canto]
Use um bloco para levar o gatinho até o canto superior direito do palco, no ponto $x=200$, $y=150$.
\end{exercicio}

\begin{solucao}
\begin{pilha}
\sbh[scEvents]{quando \bandeira\ for clicado}
\bk[scMotion]{vá para x: \entrada{200} y: \entrada{150}}
\end{pilha}
\end{solucao}

\begin{exercicio}[Meia-volta]
Faça o gatinho andar 100 passos, dar meia-volta (180 graus) e voltar ao ponto de partida.
\end{exercicio}

\begin{solucao}
Meia-volta é 180 graus. Depois de girar, andar o mesmo tanto traz o gato de volta.
\begin{pilha}
\sbh[scEvents]{quando \bandeira\ for clicado}
\bk[scMotion]{mova \entrada{100} passos}
\bk[scMotion]{gire \entrada{180} graus}
\bk[scMotion]{mova \entrada{100} passos}
\end{pilha}
\end{solucao}

\begin{exercicio}[Triângulo]
Desenhe um triângulo com o gatinho usando \bloco[scControl]{repita}. Dica: em um triângulo com lados iguais, cada giro é de $360 \div 3 = 120$ graus.
\end{exercicio}

\begin{solucao}
\begin{pilha}
\sbh[scEvents]{quando \bandeira\ for clicado}
\sbc[scControl]{repita \entrada{3} vezes}
\bk[scMotion]{mova \entrada{100} passos}
\bk[scMotion]{gire \entrada{120} graus}
\sbcend
\end{pilha}
\end{solucao}

\begin{desafio}[Sua própria figura]
Descubra quantas vezes repetir e quantos graus girar para desenhar um \textbf{hexágono} (6 lados). Dica: $360 \div 6$.
\end{desafio}

\begin{solucao}
Seis lados, então repita 6 vezes, girando $360 \div 6 = 60$ graus a cada lado.
\begin{pilha}
\sbh[scEvents]{quando \bandeira\ for clicado}
\sbc[scControl]{repita \entrada{6} vezes}
\bk[scMotion]{mova \entrada{60} passos}
\bk[scMotion]{gire \entrada{60} graus}
\sbcend
\end{pilha}
\end{solucao}
